\section{Retry pruning and failure memory}\label{sec:C18:pruning}
Tool failures create a sequence of attempted repairs that can dominate an agent's context. Some failures are transient; others reveal a wrong assumption or an invalid action. Repeating the same call can be appropriate for a rate limit and inappropriate for a permission denial. A retry policy should classify the failure before deciding whether to repeat, revise, replan, or stop.

PruneTIR distinguishes successful error resolution, a stuck repair trajectory, and repeated failures that justify temporary tool suspension~\cite{zhang2026prunetir}. It prunes resolved error traces, can resample from an earlier state, and can stop further tool use after retry conditions. Its primary reasoning evidence is concentrated in tool-integrated tasks such as mathematics. The observed correlation between erroneous calls and wrong answers does not by itself prove that deleting errors causes correctness; the interventions and ablations supply the relevant additional evidence.

Pruning has a semantic cost. A resolved error can contain a negative lesson needed to prevent recurrence. Removing the detailed trace from active context can be reasonable while preserving an explicit failure state or procedural warning outside it. Conversely, retaining every stack trace can bias the model toward a failed approach. The controller should distinguish temporary execution noise from evidence that invalidates the current plan.

Resampling from an earlier state is safe only when the environment is also restored or when the removed actions had no relevant effects. Rewinding the conversation does not rewind an external service. If a failed-looking operation partially committed, retrying an alternative can create inconsistent state. A tool-pruning method evaluated in a controlled interpreter should not automatically be generalized to irreversible workflows. State restoration and effect reconciliation belong in the deployment contract.

A proposed failure memory stores error class, attempted action, relevant state version, repair outcome, and conditions under which retry became valid. It should avoid turning an environment-specific failure into a global tool prohibition. If the tool version changes, old failures may become obsolete. This connects retry control to procedural memory and temporal validity rather than treating it as a token-removal heuristic.

Testing should retain both benign transient failures and persistent semantic failures. Compare unbounded retries, fixed retry limits, typed retry rules, and learned failure-aware control. Measure completion, wasted calls, duplicated effects, recurrence of the same error, and premature abandonment. A method that reduces calls by refusing difficult tasks needs to be distinguished from one that reaches the same success rate more efficiently. All failures and abstentions remain in the denominator.

